\documentclass[12pt]{article}
\usepackage{latexsym, amscd, %graphicx, 
amssymb}
\newcommand{\SECT}[1]
		{\vspace{.3cm}
		$\!\!\!\!\!\!\!\!\!\!\!${\bf {#1}}
		\vspace{.2cm}}

\begin{document}


\begin{center}

{\large \bf Math 152, Fall, 2004, Workshop 2\par}
\vspace{.5cm}
{{\bf Honors Section} 
\par}
\end{center}


\begin{enumerate}

\item Start with the region $\mathcal{A}$ in the first enclosed by 
the $x$-axis and the parabola $y=2x(2-x)$. Then obtain 
solids of revolution $\mathcal{S}_{1}$, $\mathcal{S}_{2}$,
and $\mathcal{S}_{3}$ by rotating $\mathcal{A}$ about the 
line 
$$y=4, \quad\quad y=-2,\quad\quad\mbox{\rm and}\quad\quad x=4,$$
respectively. All three solids are (unusual) ``doughnuts'' which
are $8$ units across, whose holes are $4$ units across, and whose 
heights are $2$ units. Sketch them.

\begin{enumerate}

\item Which do you expect to have larger volume, 
$\mathcal{S}_{1}$ or $\mathcal{S}_{2}$? Compute
their volume and check your guess.

\item Compute the volume of $\mathcal{S}_{3}$.

\end{enumerate}



\item Consider the first quadrant region $\mathcal{A}$
bounded by the curve $y =
x^{2}$, the tangent line to this curve at 
(1, 1), and the $x$-axis. 
Sketch this region.

\begin{enumerate}



\item Set up an integral giving the volume 
of the solid obtained by rotating $\mathcal{A}$ about the 
$y$-axis using the method of washers.

\item Set up an integral giving the volume 
of the solid obtained by rotating $\mathcal{A}$ about the 
$y$-axis using the method of cylindrical shells.

\item Compute one of the integrals above to find this volume.

\item Find the volume of the solid obtain by rotating $\mathcal{A}$
about $x$-axis. You may use any method.

\end{enumerate}

\end{enumerate}


\end{document}



